% Mathematical Formalization for Knowledge Graph-Based DDoS Detection
% ====================================================================
% This document contains the formal mathematical definitions, theorems,
% and proofs for the Knowledge Graph-based Layer 7 DDoS detection approach.
%
% Author: [Your Name]
% Paper: "Knowledge Graph-Based Detection and Explanation of Layer 7 DDoS Attacks"
% Target: Qualis A2 Journal (Computers & Security / J. Network and Computer Applications)

\documentclass[11pt]{article}
\usepackage{amsmath,amssymb,amsthm}
\usepackage{algorithm}
\usepackage{algorithmic}
\usepackage{booktabs}
\usepackage{tikz}
\usetikzlibrary{shapes,arrows,positioning}

% Theorem environments
\newtheorem{definition}{Definition}
\newtheorem{theorem}{Theorem}
\newtheorem{lemma}{Lemma}
\newtheorem{corollary}{Corollary}
\newtheorem{proposition}{Proposition}

\title{Mathematical Formalization:\\
Knowledge Graph-Based Detection of Layer 7 DDoS Attacks}
\author{[Author Name]}
\date{\today}

\begin{document}

\maketitle

\section{Formal Definitions}

\subsection{Knowledge Graph Definition}

\begin{definition}[Layer 7 DDoS Knowledge Graph]
A \emph{Layer 7 DDoS Knowledge Graph} is a directed labeled multigraph defined as a tuple:
\begin{equation}
    KG = (E, R, P, \mathcal{O}, \phi)
\end{equation}
where:
\begin{itemize}
    \item $E = \{e_1, e_2, \ldots, e_n\}$ is the finite set of entities (nodes)
    \item $R = \{r_1, r_2, \ldots, r_m\}$ is the finite set of relation types (edge labels)
    \item $P: E \rightarrow 2^{\mathcal{A}}$ is the property function mapping entities to attribute sets
    \item $\mathcal{O}$ is the ontology defining entity types and relation constraints
    \item $\phi: E \times R \times E \rightarrow [0,1]$ is the confidence function
\end{itemize}
\end{definition}

\subsection{Entity Types}

\begin{definition}[Entity Type Hierarchy]
The entity type set $\mathcal{T}$ is organized in a hierarchy:
\begin{equation}
    \mathcal{T} = \mathcal{T}_{HTTP} \cup \mathcal{T}_{DNS} \cup \mathcal{T}_{Session} \cup \mathcal{T}_{Attack} \cup \mathcal{T}_{Mitigation}
\end{equation}

where:
\begin{align}
    \mathcal{T}_{HTTP} &= \{\text{HTTPRequest}, \text{Endpoint}, \text{HTTPHeader}, \text{UserAgent}\} \\
    \mathcal{T}_{DNS} &= \{\text{DNSQuery}, \text{DNSDomain}, \text{DNSServer}, \text{DNSSubdomain}\} \\
    \mathcal{T}_{Session} &= \{\text{ApplicationSession}, \text{SessionToken}, \text{Identity}\} \\
    \mathcal{T}_{Attack} &= \{\text{HTTPFlood}, \text{LoginFlood}, \text{QNameRandomization}, \ldots\} \\
    \mathcal{T}_{Mitigation} &= \{\text{WAFRule}, \text{RateLimitPolicy}, \text{DNSFirewall}, \ldots\}
\end{align}
\end{definition}

\subsection{Relation Types}

\begin{definition}[Semantic Relations]
The relation set $R$ consists of semantic relations:
\begin{equation}
    R = R_{topo} \cup R_{attack} \cup R_{detect} \cup R_{mitigate}
\end{equation}

with cardinality constraints:
\begin{align}
    \text{dom}: R &\rightarrow \mathcal{T} \quad \text{(domain)} \\
    \text{range}: R &\rightarrow \mathcal{T} \quad \text{(range)}
\end{align}
\end{definition}

\section{Anomaly Detection Model}

\subsection{Anomaly Score Function}

\begin{definition}[Graph-Based Anomaly Score]
The anomaly score for an entity $e \in E$ is defined as:
\begin{equation}
    A(e) = \alpha \cdot V(e) + \beta \cdot S(e) + \gamma \cdot B(e) + \delta \cdot I(e)
\end{equation}
where:
\begin{itemize}
    \item $V(e)$: Volume-based score (traffic deviation)
    \item $S(e)$: Structural score (graph centrality deviation)
    \item $B(e)$: Behavioral score (pattern deviation)
    \item $I(e)$: Intelligence score (threat reputation)
    \item $\alpha + \beta + \gamma + \delta = 1$: Weight coefficients
\end{itemize}
\end{definition}

\subsection{Volume Score}

\begin{definition}[Volume Deviation Score]
For an endpoint entity $e$, the volume score is:
\begin{equation}
    V(e) = \sigma\left(\frac{|\mathbf{x}_e - \boldsymbol{\mu}_e|}{\boldsymbol{\sigma}_e}\right)
\end{equation}
where:
\begin{itemize}
    \item $\mathbf{x}_e$: Current feature vector for entity $e$
    \item $\boldsymbol{\mu}_e$: Historical mean vector
    \item $\boldsymbol{\sigma}_e$: Historical standard deviation
    \item $\sigma$: Sigmoid normalization function
\end{itemize}
\end{definition}

\subsection{Structural Score}

\begin{definition}[Graph Centrality Deviation]
The structural anomaly score measures deviation in graph centrality:
\begin{equation}
    S(e) = 1 - \exp\left(-\frac{|C_t(e) - C_{base}(e)|}{C_{max}}\right)
\end{equation}
where $C_t(e)$ is the centrality at time $t$ and $C_{base}(e)$ is the baseline centrality.
\end{definition}

\subsection{Behavioral Score}

\begin{definition}[Behavioral Pattern Score]
The behavioral score captures deviation from learned patterns:
\begin{equation}
    B(e) = -\log P(\text{behavior}_t | \text{behavior}_{t-1}, \ldots, \text{behavior}_{t-w})
\end{equation}
where $w$ is the observation window size.
\end{definition}

\section{Detection Rules}

\subsection{Semantic Detection Rules}

\begin{definition}[Detection Rule]
A detection rule $\rho$ is a tuple:
\begin{equation}
    \rho = (\text{pattern}, \text{threshold}, \text{attack\_type}, \text{confidence})
\end{equation}
where pattern is a graph query expressed in SPARQL-like syntax.
\end{definition}

\subsubsection{HTTP Flood Detection Rule}

\begin{equation}
    \rho_{HTTP\_Flood}: \frac{\text{count}(\text{requests}_e)}{\text{duration}_s} > \theta_{rate} \land \text{diversity}(\text{paths}) < \theta_{div}
\end{equation}

\subsubsection{Login Flood Detection Rule}

\begin{equation}
    \rho_{Login\_Flood}: \frac{\text{count}(\text{failed\_logins}_i)}{\text{time\_window}} > \theta_{login} \land \text{distinct\_IPs} > \theta_{IPs}
\end{equation}

\subsubsection{QName Randomization Detection Rule}

\begin{equation}
    \rho_{QName}: H(\text{subdomain}) > \theta_{entropy} \land \frac{|\text{unique\_subdomains}|}{|\text{queries}|} > \theta_{uniqueness}
\end{equation}

where $H(\cdot)$ denotes Shannon entropy:
\begin{equation}
    H(X) = -\sum_{x \in X} p(x) \log_2 p(x)
\end{equation}

\section{Complexity Analysis}

\begin{theorem}[Detection Complexity]
Given a knowledge graph $KG = (E, R, P, \mathcal{O}, \phi)$ with $n = |E|$ entities and $m = |R|$ relation types, the anomaly detection algorithm runs in:
\begin{equation}
    T(n, m) = O(n \cdot m \cdot k)
\end{equation}
where $k$ is the average number of relations per entity.
\end{theorem}

\begin{proof}
The detection algorithm:
\begin{enumerate}
    \item Iterates over all $n$ entities: $O(n)$
    \item For each entity, checks all incident relations: $O(m)$ average
    \item Each rule evaluation involves at most $k$ neighbors: $O(k)$
\end{enumerate}
Therefore, total complexity is $O(n \cdot m \cdot k)$.
\end{proof}

\begin{corollary}[Real-Time Detection]
For bounded average degree $k_{max}$, detection is linear in the number of entities:
\begin{equation}
    T(n) = O(n)
\end{equation}
enabling real-time detection for streaming data.
\end{corollary}

\section{Correctness Properties}

\begin{theorem}[Soundness]
If the knowledge graph $KG$ triggers detection rule $\rho$ for attack type $A$, then with probability at least $1 - \epsilon$, an attack of type $A$ is occurring:
\begin{equation}
    P(\text{Attack} = A | \rho \text{ triggered}) \geq 1 - \epsilon
\end{equation}
\end{theorem}

\begin{theorem}[Completeness]
For any attack $A$ in the attack taxonomy $\mathcal{T}_{Attack}$, there exists at least one detection rule $\rho$ such that:
\begin{equation}
    A \text{ occurring} \Rightarrow \rho \text{ triggers within } \Delta t
\end{equation}
where $\Delta t$ is the detection latency bound.
\end{theorem}

\section{Graph Construction Algorithm}

\begin{algorithm}[H]
\caption{Knowledge Graph Construction}
\begin{algorithmic}[1]
\REQUIRE Traffic logs $\mathcal{L}$, Ontology $\mathcal{O}$
\ENSURE Knowledge Graph $KG$
\STATE $KG \leftarrow \emptyset$
\FOR{each log entry $l \in \mathcal{L}$}
    \STATE $e \leftarrow \text{ExtractEntity}(l, \mathcal{O})$
    \STATE $KG.E \leftarrow KG.E \cup \{e\}$
    \STATE $R \leftarrow \text{ExtractRelations}(l, e)$
    \STATE $KG.R \leftarrow KG.R \cup R$
    \FOR{each relation $r \in R$}
        \STATE $\phi(r) \leftarrow \text{ComputeConfidence}(r)$
    \ENDFOR
\ENDFOR
\STATE $KG \leftarrow \text{EnrichWithThreatIntel}(KG)$
\RETURN $KG$
\end{algorithmic}
\end{algorithm}

\section{Detection Algorithm}

\begin{algorithm}[H]
\caption{Semantic Anomaly Detection}
\begin{algorithmic}[1]
\REQUIRE Knowledge Graph $KG$, Rules $\mathcal{R}$, Thresholds $\Theta$
\ENSURE Anomaly Set $\mathcal{A}$
\STATE $\mathcal{A} \leftarrow \emptyset$
\FOR{each entity $e \in KG.E$}
    \STATE $A(e) \leftarrow \text{ComputeAnomalyScore}(e, KG)$
    \IF{$A(e) > \theta_{anomaly}$}
        \FOR{each rule $\rho \in \mathcal{R}$}
            \IF{$\text{EvaluateRule}(\rho, e, KG)$}
                \STATE $a \leftarrow \text{CreateAnomaly}(e, \rho, A(e))$
                \STATE $\mathcal{A} \leftarrow \mathcal{A} \cup \{a\}$
            \ENDIF
        \ENDFOR
    \ENDIF
\ENDFOR
\RETURN $\mathcal{A}$
\end{algorithmic}
\end{algorithm}

\section{Explainability Model}

\begin{definition}[Explanation Graph]
For a detected anomaly $a$, the explanation graph $G_{exp}$ is:
\begin{equation}
    G_{exp}(a) = (E_{exp}, R_{exp})
\end{equation}
where:
\begin{align}
    E_{exp} &= \{e \in KG.E : \text{contributes}(e, a)\} \\
    R_{exp} &= \{r \in KG.R : \text{evidence}(r, a)\}
\end{align}
\end{definition}

\begin{definition}[Explanation Score]
The explanation score for evidence $e$ supporting anomaly $a$:
\begin{equation}
    \text{expl}(e, a) = \frac{\text{contribution}(e, a)}{\sum_{e' \in E_{exp}} \text{contribution}(e', a)}
\end{equation}
\end{definition}

\section{Performance Bounds}

\begin{theorem}[False Positive Rate Bound]
Under the assumption of independent rule triggers, the false positive rate is bounded by:
\begin{equation}
    FPR \leq \sum_{\rho \in \mathcal{R}} P(\rho \text{ false trigger}) \leq |\mathcal{R}| \cdot \epsilon_{max}
\end{equation}
where $\epsilon_{max}$ is the maximum individual rule false positive rate.
\end{theorem}

\begin{theorem}[Detection Latency]
The expected detection latency for attack $A$ is:
\begin{equation}
    E[\text{latency}_A] \leq \frac{w_{min}}{r_{attack}} + t_{process}
\end{equation}
where:
\begin{itemize}
    \item $w_{min}$: Minimum observations needed
    \item $r_{attack}$: Attack request rate
    \item $t_{process}$: Graph processing time
\end{itemize}
\end{theorem}

\section{Ontology Alignment}

\begin{definition}[STIX 2.1 Alignment]
The ontology $\mathcal{O}$ is aligned with STIX 2.1 via mapping function:
\begin{equation}
    \psi: \mathcal{T}_{Attack} \rightarrow \text{STIX}_{AttackPattern}
\end{equation}

Example mappings:
\begin{align}
    \psi(\text{HTTPFlood}) &= \text{T1498.001} \\
    \psi(\text{LoginFlood}) &= \text{T1110.001} \\
    \psi(\text{QNameRandomization}) &= \text{T1498.002}
\end{align}
\end{definition}

\section{Summary}

This formalization provides:
\begin{enumerate}
    \item Rigorous mathematical foundation for the knowledge graph model
    \item Provable bounds on detection complexity and accuracy
    \item Formal semantics for detection rules and explanations
    \item Alignment with established threat intelligence standards
\end{enumerate}

\end{document}
